\documentclass[11pt]{article}
\usepackage[T1]{fontenc}
\usepackage[utf8]{inputenc}
\usepackage{lmodern}
\usepackage[margin=1in]{geometry}
\usepackage{amsmath,amssymb,amsthm}
\usepackage{microtype}
\usepackage[hidelinks]{hyperref}
\hypersetup{pdftitle={Missing faces and topological embeddings of sphere skeleta},
  pdfauthor={Alec Kriebel},
  pdfsubject={Unrefereed research note; AI-assisted mathematical and finite verification},
  pdfkeywords={simplicial complexes, topological embeddings, van Kampen obstruction}}

\newtheorem{theorem}{Theorem}
\newtheorem{proposition}[theorem]{Proposition}
\newtheorem{lemma}[theorem]{Lemma}
\theoremstyle{remark}
\newtheorem{remark}[theorem]{Remark}
\newcommand{\F}{\mathbb F}
\newcommand{\Q}{\mathbb Q}
\newcommand{\Z}{\mathbb Z}
\newcommand{\R}{\mathbb R}
\newcommand{\lk}{\operatorname{lk}}
\newcommand{\sd}{\operatorname{sd}}
\newcommand{\interior}{\operatorname{int}}
\newcommand{\del}{\mathrm{del}}
\newcommand{\vk}{o^4_{\F_2}}
\newcommand{\supp}{\operatorname{supp}}

\title{Missing faces and topological embeddings of sphere skeleta}
\author{Alec Kriebel\\
\small Independent researcher\\
\small ORCID: \href{https://orcid.org/0009-0001-9320-500X}{0009-0001-9320-500X}}
\date{12 October 2026\\\small Version 1.0}

\begin{document}
\maketitle

\begin{abstract}
Let $S$ be a finite simplicial triangulation of the topological sphere
$S^{2d}$, and let $T$ be a missing $d$-face. We prove that adjoining $T$ to
the $d$-skeleton of $S$ prevents a topological embedding in $S^{2d}$.
For $d\geq3$, a graded Koszul calculation combines with established
ambient-completion, anisotropy and metastable embedding theorems. For
$d=2$, we apply the Nevo--Wagner bistellar transfer theorem in the PL
class supplied by the given triangulation: a contractible complementary
manifold supplies the required disk map, and a flag subdivision supplies
the base. This proves the statement of their Conjecture~1.3 without a
standard PL structure assumption on the source. No claim of first priority
or a new general method is made.
\end{abstract}

\section{Statement and relation to earlier work}

A \emph{missing $d$-face} of a simplicial complex $S$ is a subset
$T\subseteq V(S)$ with $|T|=d+1$, $T\notin S$, and every proper subset of
$T$ in $S$. Write $S^{(d)}$ for the $d$-skeleton and $|S|$ for the
realization. Embeddings below are arbitrary topological embeddings unless
the PL category is specified.

\begin{theorem}\label{thm:main}
Let $d$ be a nonnegative integer, and let $S$ be a finite simplicial complex with
$|S|\cong S^{2d}$. If $T$ is a missing $d$-face on the actual vertex set of
$S$, then
\[
                 |S^{(d)}\cup\{T\}|\not\hookrightarrow S^{2d}.
\]
\end{theorem}

In particular, no combinatorial or standard PL structure is assumed for
the higher-dimensional source. The actual-vertex convention makes the
$d=0$ assertion vacuous.

The missing-middle-face statement appears in Nevo's 2006 preprint as
Conjecture~3.8, renumbered Conjecture~3.7 in its 2007 publication
\cite{N}. Nevo--Wagner proved the conclusion for a source that is a PL
sphere \cite[Theorem~1.2]{NW}; their Conjecture~1.3 asks for the arbitrary
topological-sphere source used here. Their PL-source theorem already
forbids topological embeddings. Thus changing the category of the
prospective embedding alone does not describe the distinction at issue.

The algebraic part follows established ideas. Adiprasito's ambient
Lefschetz quotient mechanism \cite[Observation~4.3 and Corollary~4.8]{A}
already relates subcomplexes of Lefschetz spheres to graded dimension
inequalities. Adiprasito--Pat\'akov\'a \cite[Theorem~2]{AP} record an
ambient completion retaining the original abstract complex, and describe
that completion as a result sometimes regarded as folklore. Karu--Xiao
\cite[Theorems~1.2--1.3]{KX} provide the precise generic anisotropy input
used below. With Weber's existence theorem \cite[Theorem~1]{W}, these
inputs give a synthesis for $d\geq3$.

The dimension-four argument retains the Nevo--Wagner transfer and local
linking construction. We verify its hypotheses using a contractible
complement and begin at $\sd S$, within the actual PL class of $S$.
Neither the global PL standardness of a triangulated topological
4-sphere nor a standard PL complementary 4-ball is required.

\section{The graded jump and the PL obstruction}

We first work with $d\geq1$. For a field $F$, write $F[X]$ for the
Stanley--Reisner ring, with all variables of degree one. The local
homology of a triangulated topological manifold gives, for a $k$-face
$\sigma$ and $x\in\interior|\sigma|$,
\[
H_i(|S|,|S|\setminus\{x\};\Z)
 \cong\widetilde H_{i-k-1}(\lk_S\sigma;\Z).
\]
Consequently $S$ is a local integral homology sphere. This conclusion
does not identify its links as PL spheres in higher dimensions.
Reisner's criterion and the homology-sphere duality give a
Cohen--Macaulay face ring of dimension $r=2d+1$ whose Artinian reduction
has symmetric Hilbert function $h_i(S)$; in particular
$h_d(S)=h_{d+1}(S)$ (see also \cite[Sections~2.1 and 2.3]{KX}).

\begin{lemma}\label{lem:jump}
Set $K=S^{(d)}$, $L=K\cup\{T\}$, $A=F[S]$, $C=F[K]$, and $B=F[L]$.
Let $\Theta=(\theta_1,\ldots,\theta_r)$ be a regular linear parameter
sequence on $A$. Then
\begin{align}
\dim_F(B/(\Theta))_d&=h_d(S),\label{eq:degree-d}\\
\dim_F(B/(\Theta))_{d+1}&=h_{d+1}(S)+1=h_d(S)+1.\label{eq:jump}
\end{align}
The hypothesis can be fulfilled over an infinite field or a suitable
coefficient extension; it does not assert generic linear parameters over
every finite field without extension.
\end{lemma}

\begin{proof}
$A$ and $C$ have the same graded pieces and multiplication through total
degree $d+1$: taking the skeleton first removes faces with $d+2$ vertices.
$B$ and $C$ agree through degree $d$. The kernel $M$ of the surjection
$B\to C$ therefore satisfies
\[
M_j=0\quad(j\leq d),\qquad M_{d+1}=F x_T,
\quad x_T=\prod_{v\in T}x_v.
\]
The total-degree-$(d+1)$ part computing first Koszul homology is
\[
C_{d-1}\otimes\bigwedge^2 F^r
 \longrightarrow C_d\otimes F^r\longrightarrow C_{d+1}.
\]
It is identical for $A$, so regularity on $A$ yields
$H_1(\Theta;C)_{d+1}=0$. Also
$H_0(\Theta;M)_{d+1}=M_{d+1}$, since $M_d=0$.
The long exact Koszul sequence for $0\to M\to B\to C\to0$ now gives
\[
0\longrightarrow F x_T\longrightarrow(B/(\Theta))_{d+1}
 \longrightarrow(C/(\Theta))_{d+1}\longrightarrow0.
\]
The reductions of $A,C$ agree in degrees $d,d+1$, and those of $B,C$
agree in degree $d$. Their Hilbert functions give the claims.
In particular, counting an extra monomial before quotienting is not
being used to presume that it survives the quotient.
\end{proof}

\begin{proposition}\label{prop:PL}
For every $d\geq1$ under Theorem~\ref{thm:main}'s hypotheses, $L$ has
no PL embedding in the standard PL sphere $S^{2d}$.
\end{proposition}

\begin{proof}
Suppose such an embedding exists. By \cite[Theorem~2]{AP}, there is a
PL triangulation $P$ of the standard sphere containing the original
abstract $L$ as a subcomplex, without subdividing $L$. The embedding map
may change. Relabel so $V(S)=V(L)\subseteq V(P)=\{1,\ldots,N\}$ and set
\[
F=\Q(a_{ij}:1\leq i\leq r,\ 1\leq j\leq N),\qquad
\theta_i=\sum_{j=1}^N a_{ij}x_j,\qquad \ell=\sum_{j=1}^N x_j,
\]
with algebraically independent coefficients. Restriction to $V(S)$
leaves algebraically independent coefficients. On every facet of $S$,
the $r\times r$ coefficient determinant is a nonzero polynomial;
the restrictions are an l.s.o.p. on $F[S]$ and hence a regular sequence.
Extra coefficient columns merely extend the field.

The PL sphere $P$ is in particular an $\F_2$ homology sphere.
Karu--Xiao's characteristic-zero statement
\cite[Theorems~1.2--1.3]{KX}, with $n=2d+1$, $m=d$ and $n-2m=1$,
gives anisotropy of
\[
u\longmapsto\ell u^2\in(F[P]/(\Theta))_{2d+1},
\qquad u\in(F[P]/(\Theta))_d.
\]
Thus $\ell u=0$ implies $u=0$. Symmetry of the ambient Hilbert function
turns this injection into an isomorphism from degree $d$ to $d+1$.

The natural graded subcomplex quotient $F[P]\to F[L]$ sets outside
variables to zero and imposes the additional nonface relations of $L$.
It descends modulo $\Theta$. Surjectivity of multiplication by $\ell$
descends as well: lift a target class upstairs, take an upstairs
preimage, and project. Hence
\[
\ell:(F[L]/(\Theta))_d\twoheadrightarrow(F[L]/(\Theta))_{d+1},
\]
contrary to Lemma~\ref{lem:jump}. These are $r$ restricted ambient
parameters on $F[L]$; for $d\geq1$ they are overlong and are not an
l.s.o.p. on that ring, whose Krull dimension is $d+1$.
\end{proof}

\section{Topological embeddings in the metastable range}

For $d\geq1$, an embedded $d$-dimensional finite polyhedron cannot fill
$S^{2d}$, by invariance of topological dimension. Delete a point outside
the image to obtain an embedding $f:|L|\hookrightarrow\R^{2d}$.
Its normalized difference
\[
(x,y)\longmapsto\frac{f(x)-f(y)}{\|f(x)-f(y)\|}
\]
is a continuous equivariant map from the full ordered deleted product
$\{(x,y)\in|L|^2:x\ne y\}$ to $S^{2d-1}$.
Weber's original Theorem~1, printed p.~2 of \cite{W}, gives existence
of a semilinear embedding of a finite $n$-polyhedron in $\R^m$ from
this input when $2m\geq3(n+1)$. Semilinear means PL after subdivision,
as defined in his Section~1. Here
\[
n=d,\quad m=2d,\quad 4d\geq3d+3\quad\Longleftrightarrow\quad d\geq3.
\]
The equality case $d=3$ is included. A compact PL embedding in Euclidean
space lies in a PL chart of the standard sphere, contradicting
Proposition~\ref{prop:PL}. This is only an existence use of Weber's
theorem; it makes no approximation or isotopy assertion about $f$.

\section{The dimension-four obstruction}

All chains in this section have coefficients $\F_2$. For a finite
2-complex $X$, let $X^2_{\del}/\Z_2$ be the simplicial deleted product
with unordered cells $\{\sigma,\tau\}$ for vertex-disjoint simplices.
The mod-2 van Kampen class $\vk(X)$ is represented in degree four by
the cochain
\[
\varphi_f\{\sigma,\tau\}=f(\sigma)\cdot f(\tau)\pmod2
\]
for a general-position map into $\R^4$. Its class is independent of
the representative. A topological embedding forces this class to
vanish; see \cite[Section~2]{NW}. We will prove the stronger assertion
$\vk(S^{(2)}\cup\{M\})\ne0$ for every missing triangle $M$ of $S$.

\subsection{The source structure and a complement lemma}

A finite triangulation of a topological 4-manifold is a combinatorial
4-manifold \cite[p.~797]{DFL}. To recall the dimension-specific point,
positive-dimensional face links have sphere homology in dimensions at
most two. Inductively they are PL spheres; vertex links are therefore
closed triangulated homology 3-spheres. At a vertex, choose nested
coordinate and conical neighborhoods. Place a link loop radially in a
small cone contained in a punctured coordinate 4-ball, itself contained
in the punctured open star. The coordinate ball makes the loop
nullhomotopic, and radial retraction puts the nullhomotopy in the link.
The link is simply connected. Three-dimensional Poincar\'e and PL
uniqueness then make it a PL 3-sphere. This supplies a PL structure on
$S$ without identifying its global PL class with the standard sphere.

\begin{lemma}\label{lem:complement}
Let $Q$ be a combinatorial 4-manifold with $|Q|\cong S^4$, and let
$B_0=\alpha*\partial\beta$ be the closed support of a legal bistellar
move. Then $C=\overline{|Q|\setminus\interior|B_0|}$ is a contractible
compact PL manifold with boundary and collar.
\end{lemma}

\begin{proof}
The support is a standard PL 4-ball. Its boundary is locally flat:
at a boundary vertex its link pair is a PL 2-sphere in a PL 3-sphere,
where three-dimensional PL Schoenflies applies; higher-dimensional
strata reduce to the lower-dimensional link pairs. Thus $C$ is a
compact PL 4-manifold with boundary $S^3$ and a collar.
The collared decomposition $|Q|=|B_0|\cup C$ gives connectedness and
$\pi_1(C)=0$ by Mayer--Vietoris and van Kampen. The fundamental-class
map
\[
H_4(|Q|;\Z)\longrightarrow H_4(|Q|,C;\Z)
 \cong H_4(|B_0|,\partial|B_0|;\Z)
\]
has degree $\pm1$. The pair sequence therefore gives
$\widetilde H_*(C;\Z)=0$, including degree three. For a simply connected
CW complex, acyclicity and Hurewicz give vanishing of every homotopy
group; Whitehead gives contractibility.
\end{proof}

Only contractibility is claimed. Local three-dimensional Schoenflies
has not been turned into a global PL four-dimensional Schoenflies
assertion.

\subsection{Common missing triangles and the transfer hypotheses}

Define $E(Q)$ to mean that every missing triangle $M$ of $Q$ satisfies
$\vk(Q^{(2)}\cup\{M\})\ne0$. Let a move replace
\[
B_0=\alpha*\partial\beta\quad\hbox{by}\quad
B_1=\partial\alpha*\beta,\qquad
\dim\alpha=p,\quad\dim\beta=q,\quad p+q=4,
\]
and denote the result by $Q'$. First suppose $M$ is missing in both.
Set
\[
\mathcal K=Q^{(2)}\cup\{M\},\quad
\mathcal L=(Q')^{(2)}\cup\{M\},\quad
R=B_0^{(2)},\quad R'=B_1^{(2)}.
\]
Theorem~5.3 of \cite{NW} applies to arbitrary finite 2-complexes with
this skeletal move, provided its structural setup and the following
conditions hold: $R$ is induced in $\mathcal K$; the open star of
$\alpha$ in $\mathcal K$ lies in $R$; general-position maps $f,g$ agree
on the common complex and have zero intersection cochain on disjoint
triangle pairs involving $R,R'$ respectively; an outside vertex exists;
and
\begin{equation}\label{eq:homology}
\widetilde H_i(\mathcal K-W;\F_2)=0\quad(i=0,1),
\qquad W=R_0\setminus\rho,\quad\rho\in R\cap R'.
\end{equation}
Here deleting $W$ means taking the induced complex on the remaining
vertices. We check every condition.

\paragraph{Inducedness and the removed star.}
Before the move $\beta$ is a missing face (or a new vertex in a vertex
insertion). On the old support vertices, every subset not containing
all of $\beta$ belongs to $B_0$, so the old support is induced.
A common missing triangle contained on those vertices would, by
minimality, have to equal $\beta$, which becomes a face. Thus adjoining
$M$ preserves inducedness. A common $M$ cannot contain $\alpha$:
equality contradicts that $\alpha$ was an old face, and proper
containment would give a proper face of $M$ that disappears. The legal
move's open-star condition therefore persists after augmentation.

\paragraph{The maps.}
All proper faces of $M$ survive, so $|\partial M|$ lies in the
complement $C$ of Lemma~\ref{lem:complement}. Contractibility extends
its prescribed inclusion to a disk \emph{map}. A collar pushes the
disk interior into $\interior C$, and relative PL general position
makes that interior avoid the old 1-skeleton, since $2+1<4$. Its
boundary stays fixed in the 1-skeleton; disk self-intersections are
permitted. Choose $x\in\interior|B_0|$ off the old 2-skeleton and use
a topological homeomorphism $|Q|\setminus\{x\}\cong\R^4$. Map the old
skeleton naturally and $M$ by this disk. No vertex-disjoint triangle
pair involving $R$ intersects.

A homeomorphism $|B_1|\to|B_0|$ fixed on their common boundary defines
the new support map; use the old map on the common exterior and $M$.
The maps agree on $\mathcal K\cap\mathcal L$, and the corresponding
disjoint pairs involving $R'$ also have disjoint images. To get
general position, regard these as one continuous map on the finite
abstract union $\mathcal K\cup\mathcal L$ and approximate it
simultaneously. Each forbidden pair has disjoint compact images,
hence positive separation, and there are finitely many pairs. A
sufficiently close general-position PL approximation preserves all
separations and agreement upon restriction. The potentially non-PL
puncturing homeomorphism need not be kept fixed by that approximation.

\paragraph{Outside vertices, homology and cones.}
An induced support ball cannot exhaust the vertices of a closed
sphere, since inducedness would then identify the sphere with the
ball. The outside vertex sets before and after the move agree,
including vertex insertion and deletion.
For $\rho\in R\cap R'\subseteq(\partial\alpha*\partial\beta)^{(2)}$
and $W=R_0\setminus\rho$, the induced old complex is
\[
Q[W]=(\alpha\setminus\rho)*(\partial\beta\setminus\rho).
\]
The first factor is a nonempty simplex: a common-boundary face cannot
contain all of $\alpha$. Thus $Q[W]$ is a cone. Alexander duality in
the topological sphere makes $|Q|\setminus|Q[W]|$ acyclic.
It deformation retracts onto $|Q-W|$: in barycentric coordinates,
retain coordinates outside $W$ and normalize their positive sum;
the straight homotopy stays in the complement. The 2-skeleton has
the same $H_0,H_1$, and adding a triangle with its boundary present
preserves $H_0$ and can only quotient $H_1$. This proves
\eqref{eq:homology}.

The same argument works in the old \emph{unaugmented} skeleton. It
also supplies the precise cones of \cite[Lemma~5.2]{NW}: for an
outside vertex $v$ and common vertex or edge $\rho$, choose chains
$v\bullet\rho$ with
\begin{equation}\label{eq:cone}
\partial(v\bullet\rho)=\rho+
 \sum_{a\in\rho}v\bullet(\rho\setminus\{a\}),\qquad
\supp(v\bullet\rho)\cap(R_0\cap R'_0)\subseteq\rho.
\end{equation}
Here $v\bullet\varnothing=v$. For a vertex the boundary is $a+v$;
use an actual simplex whenever
the cone simplex exists. Otherwise the induction's boundary cycle
bounds in the relevant deleted complex by $H_0$ or $H_1$ vanishing.
For a face meeting the exterior, delete the same support vertices
using its intersection with the support. This is exactly the
support-avoidance construction, not an assumption that all geometric
cones are present.

The structural conditions and \cite[conditions (I.a)--(II.b)]{NW} are
now verified. Their Theorem~5.3 gives
$\vk(\mathcal K)\ne0\Rightarrow\vk(\mathcal L)\ne0$ for common $M$.

\subsection{All newly missing triangles and the explicit witness}

The remaining cases can be classified without a sphere-structure
assumption. If a triangle was an old face and becomes missing, it
contains $\alpha$, while every proper face survives. It must be
$M=\alpha$, with $(p,q)=(2,2)$. Otherwise some proper edge of $M$ is
new and contains $\beta$. If $q=1$ this gives $(p,q)=(3,1)$ and
$M=v*\beta$, with $v$ outside the support; a triangle wholly on the
support vertices would be a new face. If $q=0$, the new vertex of
$\beta$ is adjacent only to support vertices; any proposed missing
triangle having all three edges is a new support face. For completeness,
$(p,q)=(0,4)$ removes a vertex and $(1,3)$ removes an edge, so a
triangle involving the removed simplex cannot have all proper faces
present; no new proper edge arises in either case. These exhaust
all five move types.

When $(p,q)=(2,2)$, the old missing triangle is $\beta$ and
\[
Q^{(2)}\cup\{\beta\}=(Q')^{(2)}\cup\{\alpha\}.
\]
The induction hypothesis applies. For $(p,q)=(3,1)$ we give the
homological Flores witness explicitly, following the local mechanism
of \cite[proof of Theorem~1.2]{NW}.

Put $U=\alpha_0\cup\beta_0$ and $V=U\cup\{v\}$, so $|U|=6$ and
$|V|=7$. Every triangle on $U$ is in the new support; every edge
on $U$ except $\beta$ is in the old skeleton. The old 2-skeleton
is a subcomplex of the new one. Use \eqref{eq:cone} in that old
\emph{unaugmented} skeleton. Let $F_V$ be the formal full 2-skeleton
on $V$. Define a chain map $B$ through dimension two by fixing
vertices, support edges and support triangles, and setting
\[
B(va)=v\bullet a,\qquad
B(ve)=v\bullet e\ (e\ne\beta),\qquad B(v\beta)=M.
\]
Here $a\in U$ and $e$ is an edge on $U$; $\beta$ itself is fixed
as a new edge. Formula~\eqref{eq:cone} gives the chain identities
except possibly on $v\beta$. Both edges from $v$ to the endpoints
of $\beta$ are old edges, since they are proper faces of $M$ and
the move introduces no edges incident to $v$. The prescribed vertex
cones are those edges, so $\partial B(v\beta)=B(\partial(v\beta))$
as well.

The formal deleted-product chain
\[
\Omega=\sum_{\substack{\sigma,\tau\subset V,\ |\sigma|=|\tau|=3\\
                         \sigma\cap\tau=\varnothing}}
                    \{\sigma,\tau\},
\]
where each unordered pair is counted once, is a 4-cycle. Each boundary
cell is an edge paired with a disjoint triangle; its five vertices
leave two possible extensions of the edge to a triangle, so its
coefficient is zero modulo two. Apply $B$ to both factors to get
$\omega$. At most one factor of a disjoint pair contains $v$;
the other is supported on $U$. Support avoidance in
\eqref{eq:cone} keeps all replaced pairs vertex-disjoint, including
boundary pairs. Hence the chain map on these deleted-product cells
is defined and $\partial\omega=0$.

No cone chain contains $M$, because all cones lie in the old
unaugmented skeleton. Thus $M$ has its sole occurrence from the
explicit assignment $B(v\beta)=M$. Its partner chain consists of
the four triangles on $V\setminus M=\alpha_0$, namely
$\partial\alpha$. There is no hidden use of an old missing triangle.

Choose a puncture outside the new ball and the new 2-skeleton and
identify the punctured topological sphere with $\R^4$. Its natural
map embeds the unaugmented new skeleton. In the explicit ball
$B_1=\partial\alpha*\beta$, realize the tetrahedron $\alpha$ and
interval $\beta$ in complementary 3- and 1-dimensional subspaces,
with one transverse interior intersection. This supplies a spanning
3-disk for $\partial\alpha$ inside the ball. Pull it back to the
actual move ball. The other two edges of $\partial M$ go to the
outside vertex $v$ and miss the ball interior; their endpoints on
$\beta$ are away from $\partial\alpha$. The spanning disk therefore
meets $\partial M$ once modulo two. Consequently
\[
\lk_2(\partial M,\partial\alpha)=1.
\]
The linking/intersection boundary relation gives intersection number
one for a general-position filling of $M$ against $\partial\alpha$.
All other pairs in $\omega$ lie in the embedded unaugmented skeleton
and contribute zero. A close simultaneous general-position choice
preserves their finite compact separations and the linking number.
It follows that
\[
\langle\vk((Q')^{(2)}\cup\{M\}),\omega\rangle
 =\varphi_f(\omega)=1.
\]
This verifies the newly missing triangle case and proves
$E(Q)\Rightarrow E(Q')$ for every bistellar move.

\subsection{The base in the source's own PL class}

The barycentric subdivision $\sd S$ is flag: its vertices are faces
of $S$, and pairwise adjacency means pairwise comparability, hence a
chain. It has no missing triangles, so $E(\sd S)$ holds vacuously.
The subdivision and $S$ are closed combinatorial 4-manifolds that
are PL homeomorphic for the PL structure supplied by this
triangulation. Pachner's theorem, as stated and proved in
\cite[Theorem~5.9]{L}, gives a finite bistellar sequence between them,
allowing simplicial isomorphisms. Every intermediate complex remains
homeomorphic to $S^4$ in the same PL class. Transfer at every step
gives $E(S)$, completing $d=2$. No comparison with the standard
PL 4-sphere occurs in this induction.

\section{Low dimensions and verification scope}

For $d=1$, a triangulated 2-sphere with $n$ vertices has $3n-6$ edges,
by Euler's formula and the triangle-edge incidence count. Adding a
missing edge gives $3n-5$, exceeding the simple planar graph bound.
Thus the augmented graph does not embed in $S^2$. For $d=0$ there
is no missing singleton on the actual vertex set. If a convention
permits a ghost singleton, its addition gives three points, which
cannot embed in the two-point sphere $S^0$. These cases, the
metastable argument and the dimension-four obstruction prove
Theorem~\ref{thm:main}.

The research record includes exact finite computations of face-ring
reductions and homological witnesses, together with adverse controls.
They provide finite consistency and reproducibility evidence. The
universal conclusion rests on the arguments and cited theorems above;
those computations neither replace them nor constitute formal proof
certification.

\paragraph{Disclosure and manuscript status.}
AI tools were used extensively in solving, drafting, separate
adversarial auditing, and finite verification. This manuscript is
unrefereed, and no independent human peer review has been obtained.
AI review is not formal verification. The author is responsible for
the stated claims and citations. The note makes no assertion that a
current literature search certifies the problem's open status or that
this is its first resolution.

\clearpage
\begin{thebibliography}{9}

\bibitem{A}
K. Adiprasito,
\emph{Combinatorial Lefschetz theorems beyond positivity},
arXiv:1812.10454v4, 2 July 2019.
\url{https://arxiv.org/abs/1812.10454v4}.

\bibitem{AP}
K. Adiprasito and Z. Pat\'akov\'a,
\emph{A higher dimensional version of F\'ary's theorem},
Bull. London Math. Soc. \textbf{57} (2025), no.~5, 1409--1414.
\href{https://doi.org/10.1112/blms.70036}{doi:10.1112/blms.70036}.
Author version: \url{https://arxiv.org/abs/2404.12265}.

\bibitem{DFL}
M. W. Davis, J. Fowler and J.-F. Lafont,
\emph{Aspherical manifolds that cannot be triangulated},
Algebr. Geom. Topol. \textbf{14} (2014), no.~2, 795--803.
\href{https://doi.org/10.2140/agt.2014.14.795}{doi:10.2140/agt.2014.14.795}.

\bibitem{KX}
K. Karu and E. Xiao,
\emph{On the anisotropy theorem of Papadakis and Petrotou},
Algebraic Combinatorics \textbf{6} (2023), no.~5, 1313--1330.
\href{https://doi.org/10.5802/alco.298}{doi:10.5802/alco.298}.

\bibitem{L}
W. B. R. Lickorish,
\emph{Simplicial moves on complexes and manifolds},
Geom. Topol. Monogr. \textbf{2} (1999), 299--320.
\url{https://arxiv.org/abs/math/9911256}.

\bibitem{N}
E. Nevo,
\emph{Higher minors and Van Kampen's obstruction},
Math. Scand. \textbf{101} (2007), 161--176.
Original preprint, 23 February 2006:
\url{https://arxiv.org/abs/math/0602531v1}.

\bibitem{NW}
E. Nevo and U. Wagner,
\emph{On the embeddability of skeleta of spheres},
Israel J. Math. \textbf{174} (2009), 381--402.
\href{https://doi.org/10.1007/s11856-009-0119-5}{doi:10.1007/s11856-009-0119-5}.
The section and theorem locators used here also refer to the author
preprint dated 4 February 2008:
\url{https://pi.math.cornell.edu/~eranevo/homepage/NevoWagner-EmbeddabilitySkeletaSpheres-Rev2.pdf}.

\bibitem{W}
C. Weber,
\emph{Plongements de poly\`edres dans le domaine m\'etastable},
Comment. Math. Helv. \textbf{42} (1967), 1--27.
\href{https://doi.org/10.1007/BF02564408}{doi:10.1007/BF02564408}.

\end{thebibliography}
\end{document}
